% Starter article in the Springer LNCS layout used for the class proceedings.
% Easiest: on Overleaf, create a project from the "Springer LNCS" template
% (it provides llncs.cls and splncs04.bst), then replace its main .tex with this file.
\documentclass[runningheads]{llncs}
\usepackage[T1]{fontenc}
\usepackage{graphicx}
\usepackage{amsmath,amssymb}
\usepackage{booktabs}
\usepackage{hyperref}
\usepackage{xurl}

\begin{document}

\title{Informative Title: Method + Task + Place}
\titlerunning{Short Running Title}
\author{Your Full Name\inst{1}}
\authorrunning{Initials. Surname}
\institute{Department of Computer Science, Ramanujan College, University of Delhi, New Delhi, India\\
\email{your.email@example.com}}

\maketitle

\begin{abstract}
150--250 words: the problem and why it matters, the data, the methods compared,
the key numbers (with the baseline), and the main conclusion. No citations.
\keywords{keyword one \and keyword two \and machine learning \and SDG~X}
\end{abstract}

\section{Introduction}
Why the problem matters (with cited facts \cite{ref1}), the gap, your research question,
and your contributions as a short list. End with an outline of the paper.

\section{Related Work}
8 or more recent research papers: what they did, data, results and limitations, and how your study differs.

\section{Data and Preprocessing}
Source, link, licence and access date \cite{dataset}. Variables table, EDA, missing values,
outliers, scaling, encoding, and feature selection or PCA, with reasons.

\section{Methodology}
A workflow figure and each algorithm with its key equation, for example
\begin{equation}
  \hat{y} = \beta_0 + \sum_{j=1}^{p} \beta_j x_j .
  \label{eq:linear}
\end{equation}
Define every symbol. State the metrics.

\section{Experimental Setup}
Split and cross-validation scheme (time-ordered or grouped where needed; test set used once),
hyper-parameter grids, software versions, and random seeds.

\section{Results}
\begin{table}[htbp]
\centering
\caption{Test results (mean $\pm$ sd over 5 folds). Best value in bold.}
\label{tab:results}
\begin{tabular}{lcc}
\toprule
Model & RMSE & $R^2$ \\
\midrule
Baseline (mean) & 0.00 $\pm$ 0.00 & 0.00 \\
Model A         & 0.00 $\pm$ 0.00 & 0.00 \\
\textbf{Model B} & \textbf{0.00 $\pm$ 0.00} & \textbf{0.00} \\
\bottomrule
\end{tabular}
\end{table}
Refer to every table and figure in the text (Table~\ref{tab:results}). Include error or subgroup analysis.

\section{Discussion}
Interpretation, limitations, ethics and bias, and who could use the results.

\section{Conclusion and Future Work}
Answer the research question and give 2--3 concrete next steps.

\begin{credits}
\subsubsection{\discintname}
The author has no competing interests to declare.

\subsubsection{Data and code availability.} \url{https://github.com/your-username/your-repo}

\subsubsection{Use of AI tools.} State honestly which AI tools were used and for what,
and that you checked all output and take responsibility for the content.
\end{credits}

\begin{thebibliography}{8}
\bibitem{ref1} A. Author and B. Author, ``Title of the paper,'' \emph{Journal Name}, vol.~1, no.~1, pp.~1--10, 2024, doi:~10.xxxx/xxxxx.
\bibitem{dataset} Organisation, ``Dataset title,'' Portal name. [Online]. Available: \url{https://example.org} (accessed DD Mon YYYY).
\end{thebibliography}

\end{document}
